% Student-facing ethics pack — information sheet + informed consent (ES)
% Build: make student  (XeLaTeX — CJK for 道)
\documentclass[11pt,a4paper]{scrartcl}

\usepackage[spanish]{babel}
\usepackage{amssymb}
\usepackage{microtype}
\usepackage{geometry}
\usepackage{booktabs}
\usepackage{setspace}

% Shared chrome for CEI / student ethics PDFs — UDIT · ECSIT
% Built with XeLaTeX (CJK for 道). Input from documents in this directory.
\usepackage{graphicx}
\usepackage{fontspec}
\usepackage{hyperref}
\usepackage{xcolor}
\usepackage{tabularx}
\usepackage{fancyhdr}
\usepackage{calc}

% Single CJK glyph (道) via fontspec — no xeCJK dependency
\IfFontExistsTF{Songti SC}{%
  \newfontfamily\TTODdao{Songti SC}%
}{%
  \IfFontExistsTF{PingFang SC}{%
    \newfontfamily\TTODdao{PingFang SC}%
  }{%
    \IfFontExistsTF{STSong}{%
      \newfontfamily\TTODdao{STSong}%
    }{%
      \newcommand{\TTODdao}{}%
    }%
  }%
}

\definecolor{uditnavy}{HTML}{0B1F33}
\definecolor{uditink}{HTML}{1A1A1A}
\definecolor{rulegray}{HTML}{CCCCCC}

\newcommand{\assetspath}{../assets/}
\newcommand{\EcsitUrl}{https://www.udit.es/lineas-de-investigacion/grupo-de-investigacion-innovacion-y-tecnologia-desde-y-para-la-educacion-la-cultura-y-la-sociedad/}
\newcommand{\EcsitLink}{\href{\EcsitUrl}{udit.es $\rightarrow$ grupo ECSIT}}
\newcommand{\TTODfull}{{\TTODdao 道} The Tao of Development (TTOD)}

\newcommand{\LogoRow}{%
  \noindent
  \begin{tabularx}{\textwidth}{@{}lXr@{}}
    \includegraphics[height=0.95cm]{\assetspath udit-logo-print.png}
    & &
    \includegraphics[height=1.45cm]{\assetspath ecsit-logo-print.png}
  \end{tabularx}%
  \par\vspace{0.55em}
  \noindent{\color{rulegray}\hrule height 0.35pt}%
  \par\vspace{0.65em}
}

\newcommand{\LegalEntities}{%
  \noindent{\scriptsize\color{uditink}
  UDIT ESTUDIOS SUPERIORES INTERNACIONALES, S.L.\ · NIF B84014265 ·
  UDIT Escuela de Formación Profesional, S.L.\ · NIF B09963117.}%
}

% Role (left) | Name + email only (right) — no role text repeated
\newcommand{\AuthoritiesBlock}{%
  \noindent{\sffamily\bfseries\small Autoridades de investigación / Research authorities}\\[0.25em]
  \noindent
  \begin{tabularx}{\textwidth}{@{}>{\raggedright\arraybackslash}p{4.6cm}X@{}}
    \textbf{Coordinación grados tecnológicos} &
      Sandra Garrido ---
      \href{mailto:sandra.garrido@udit.es}{sandra.garrido@udit.es}\\[0.2em]
    \textbf{Dirección de grados} &
      Fernando Blázquez Piñeiro ---
      \href{mailto:fernando.blazquez@udit.es}{fernando.blazquez@udit.es}\\[0.2em]
    \textbf{IP del grupo ECSIT} &
      Rafael Conde Melguizo, PhD ---
      \href{mailto:rafael.conde@udit.es}{rafael.conde@udit.es}\\[0.2em]
    \textbf{CEI UDIT} &
      \href{mailto:comite.etica.investigacion@udit.es}{comite.etica.investigacion@udit.es}
  \end{tabularx}%
  \par\vspace{0.35em}%
  \LegalEntities
}

\newcommand{\StudyMetaPI}{%
  \noindent
  \begin{tabularx}{\textwidth}{@{}>{\raggedright\arraybackslash}p{4.6cm}X@{}}
    \textbf{Investigador principal} &
      Rubén Vega Balbás, PhD ---
      \href{mailto:ruben.vega@udit.es}{ruben.vega@udit.es}\\[0.15em]
    \textbf{Afiliación académica} &
      Grado Oficial en Desarrollo Full-Stack
      (\emph{La evolución lógica de la ingeniería informática hacia una
      titulación enfocada en el mundo real.})\\[0.15em]
    \textbf{Proyecto / producto} &
      \TTODfull\ · plataforma Oráculo · Cohorte Front-End~II\\[0.15em]
    \textbf{Grupo de investigación} &
      ECSIT · \EcsitLink\\[0.15em]
    \textbf{Comité} &
      \href{mailto:comite.etica.investigacion@udit.es}{comite.etica.investigacion@udit.es}\\[0.15em]
    \textbf{Código de estudio} &
      TTOD-FEII-2026-27\\[0.15em]
    \textbf{Versión del paquete} &
      2026-09-26 · cohorte pedagógica $n=8$
      (uso investigador: solo con consentimiento, $n\leq 8$)
  \end{tabularx}%
}

% Full Chicago (author–date) methods backbone — from docs/public/research/methodology.md
\newcommand{\ChicagoMethodsRefs}{%
\begin{itemize}
  \item Brown, Neil C.~C., and Mark Guzdial. 2024.
        ``Confidence vs Insight: Big and Rich Data in Computing Education Research.''
        In \emph{Proceedings of the 55th ACM Technical Symposium on Computer Science Education V.~1},
        158--64. \url{https://doi.org/10.1145/3626252.3630813}.
  \item Fischer, Björn, Berit Barthelmes, Sven Eric Panitz, Eva-Maria Iwer, and Ralf Dörner. 2024.
        ``Seeking Consent for Programming Process Data Collection with Trustee-Based Encryption.''
        In \emph{Proceedings of the 2024 ACM Conference on International Computing Education Research V.1},
        131--42. \url{https://doi.org/10.1145/3632620.3671125}.
  \item Heckman, Sarah, Jeffrey C.~Carver, Mark Sherriff, and Ahmed Al-Zubidy. 2022.
        ``A Systematic Literature Review of Empiricism and Norms of Reporting in Computing Education Research Literature.''
        \emph{ACM Transactions on Computing Education} 22 (1): Article~3.
        \url{https://doi.org/10.1145/3470652}.
  \item McGill, Monica M., Sarah Heckman, Christos Chytas, et~al. 2023.
        ``Conducting Sound, Equity-Enabling Computing Education Research.''
        In \emph{Proceedings of the 2023 Working Group Reports on Innovation and Technology in Computer Science Education},
        30--56. \url{https://doi.org/10.1145/3623762.3633495}.
  \item Runeson, Per, and Martin Höst. 2009.
        ``Guidelines for Conducting and Reporting Case Study Research in Software Engineering.''
        \emph{Empirical Software Engineering} 14 (2): 131--64.
        \url{https://doi.org/10.1007/s10664-008-9102-8}.
  \item Runeson, Per, Martin Höst, Austen Rainer, and Björn Regnell. 2012.
        \emph{Case Study Research in Software Engineering: Guidelines and Examples}.
        Hoboken, NJ: Wiley. \url{https://doi.org/10.1002/9781118181034}.
  \item Wohlin, Claes, Per Runeson, Martin Höst, Magnus C.~Ohlsson, Björn Regnell, and Anders Wesslén. 2012.
        \emph{Experimentation in Software Engineering}. Berlin: Springer.
        \url{https://doi.org/10.1007/978-3-642-29044-2}.
\end{itemize}%
}


\geometry{margin=2.2cm, top=2.0cm, bottom=2.2cm}
\setkomafont{section}{\sffamily\bfseries\large}
\setkomafont{subsection}{\sffamily\bfseries\normalsize}
\setkomafont{disposition}{\sffamily}
\setlength{\parskip}{0.35em}
\setlength{\parindent}{0pt}
\setlength{\leftmargini}{1.4em}
\setlength{\itemsep}{0.15em}
\setlength{\topsep}{0.2em}
\setstretch{1.05}

\hypersetup{
  colorlinks=true,
  urlcolor=uditnavy,
  linkcolor=uditnavy,
  pdftitle={Hoja de información y consentimiento — 道 The Tao of Development},
  pdfauthor={Rubén Vega Balbás, PhD · UDIT / ECSIT},
  pdfsubject={CEI UDIT — participantes estudiantes}
}

\pagestyle{fancy}
\fancyhf{}
\renewcommand{\headrulewidth}{0pt}
\fancyfoot[L]{\scriptsize\color{uditnavy} \TTODfull\ · CEI · ECSIT}
\fancyfoot[C]{\scriptsize\color{uditink}\thepage}
\fancyfoot[R]{\scriptsize\color{uditink}v.\ 2026-09-26}

\begin{document}

\LogoRow

\begin{center}
  {\Large\sffamily\bfseries Hoja de información a participantes}\\[0.35em]
  {\normalsize\sffamily y formulario de consentimiento informado}\\[0.6em]
  {\large\sffamily Aprendizaje haciendo --- grupo \TTODfull\ / Front-End~II}\\[0.25em]
  {\small\color{uditnavy} Estudio de caso educativo / investigación basada en diseño}
\end{center}

\vspace{0.4em}
\noindent
\begin{tabularx}{\textwidth}{@{}>{\raggedright\arraybackslash}p{4.6cm}X@{}}
  \textbf{Investigador principal} &
    Rubén Vega Balbás, PhD ---
    \href{mailto:ruben.vega@udit.es}{ruben.vega@udit.es}\\[0.12em]
  \textbf{Afiliación académica} &
    Grado Oficial en Desarrollo Full-Stack
    (\emph{La evolución lógica de la ingeniería informática hacia una
    titulación enfocada en el mundo real.})\\[0.12em]
  \textbf{Grupo de investigación} &
    ECSIT · \EcsitLink\\[0.12em]
  \textbf{Comité} &
    \href{mailto:comite.etica.investigacion@udit.es}{comite.etica.investigacion@udit.es}\\[0.12em]
  \textbf{Versión del documento} &
    2026-09-26 · cohorte pedagógica $n=8$
    (uso investigador: solo con consentimiento, $n\leq 8$)
\end{tabularx}

\vspace{0.55em}
\AuthoritiesBlock

\vspace{0.45em}
\noindent{\color{rulegray}\hrule height 0.4pt}

\section*{1.\ ¿De qué trata este estudio?}

Queremos comprender \textbf{cómo se aprende desarrollando de verdad}:
participando como desarrollador/a en una aplicación real
(la plataforma \TTODfull\ / Oráculo),
con asistencia de herramientas de IA agentica, y en grupo con otros
desarrolladores. No buscamos demostrar que <<la IA mejora las notas>>.
Buscamos describir el proceso de aprendizaje y mejorar el diseño docente.

Hay otra línea de investigación \textbf{filosófica} en el proyecto TTOD;
\textbf{no forma parte} de este estudio y \textbf{no usa} vuestro trabajo
de clase.

\section*{2.\ ¿Por qué me piden participar?}

Porque ya formáis parte de la actividad docente de Front-End~II.
La asignatura sigue igual \textbf{aceptéis o no} el uso investigador
de vuestros artefactos ni la autorización de imagen/comunicación.

\section*{3.\ ¿Qué se os pide?}

\textbf{Como estudiantes (docencia):} desarrollar, revisar, declarar el uso
de IA y defender decisiones --- lo ya previsto en la guía~/ encargo.

\textbf{Como participantes en investigación (opcional y separable):}
permitir que, \textbf{una vez autorizado el protocolo} y
\textbf{tras cerrar las calificaciones relevantes}, se analicen de forma
\textbf{pseudonimizada} (sin vuestro nombre ni identificadores directos)
artefactos de proceso ya producidos para la evaluación: commits/PRs,
declaraciones de uso de IA, notas estructuradas de defensa, y extractos
técnicos no identificativos.

\textbf{Autorización de imagen y comunicación (opcional y separable, ítem~G):}
consentir que la universidad / el proyecto puedan usar \textbf{vuestra imagen}
(fotografías o capturas de momentos del grupo) en comunicaciones
relacionadas con el proyecto: LinkedIn y otras redes, medios, publicaciones
académicas o institucionales, revistas internas de UDIT y el sitio web público,
incluyendo material gestionado por el departamento de comunicación.

En la versión mínima de este estudio \textbf{no} se piden:
\begin{itemize}
  \item encuestas sobre vuestra vida privada;
  \item datos de salud u otras categorías especiales;
  \item grabación de audio/vídeo de defensa o entrevista (salvo consentimiento
        adicional específico distinto del ítem~G).
\end{itemize}

\section*{4.\ Datos personales}

El estudio está diseñado para que \textbf{vuestros datos personales no estén
en juego} como objeto de la investigación: no se publicarán nombres, correos
ni identificadores en el corpus investigador; ese corpus, si se autoriza, estará
\textbf{pseudonimizado} y con acceso restringido. Las notas y el expediente
académico permanecen en los sistemas docentes habituales y \textbf{fuera}
del corpus de investigación.

La autorización de \textbf{imagen} (ítem~G) es distinta: ahí sí puede aparecer
vuestra imagen reconocible en canales de comunicación del proyecto/UDIT, solo
si la marcáis. Podéis aceptar B sin aceptar G, y viceversa.

Podéis retirar el consentimiento investigador o la autorización de imagen sin
efecto sobre la calificación (la retirada de imagen no obliga a retirar
piezas ya publicadas de forma irreversible en terceros, y así se os informará).

\section*{5.\ Beneficios y riesgos}

\begin{itemize}
  \item \textbf{Beneficio directo:} ninguno académico por consentir.
  \item \textbf{Beneficio colectivo:} mejorar la pedagogía; visibilidad del
        proyecto y del grupo en canales institucionales (si aceptáis G).
  \item \textbf{Riesgos:} reidentificación residual en el corpus (mitigada por
        seudonimización); incomodidad jerárquica (mitigada porque la nota no
        depende del consentimiento); exposición pública de imagen si aceptáis~G.
\end{itemize}

\section*{6.\ Voluntariedad}

La participación investigadora y la autorización de imagen son
\textbf{voluntarias}. Negaros o retiraros \textbf{no afecta} a la nota
ni al acceso a la asignatura.

\section*{7.\ Contacto}

\begin{itemize}
  \item Estudio: Rubén Vega Balbás, PhD ---
        \href{mailto:ruben.vega@udit.es}{ruben.vega@udit.es}
  \item Coordinación grados tecnológicos: Sandra Garrido ---
        \href{mailto:sandra.garrido@udit.es}{sandra.garrido@udit.es}
  \item Dirección de grados: Fernando Blázquez Piñeiro ---
        \href{mailto:fernando.blazquez@udit.es}{fernando.blazquez@udit.es}
  \item IP ECSIT: Rafael Conde Melguizo, PhD ---
        \href{mailto:rafael.conde@udit.es}{rafael.conde@udit.es}
  \item CEI:
        \href{mailto:comite.etica.investigacion@udit.es}{comite.etica.investigacion@udit.es}
  \item DPD UDIT:
        \href{mailto:dpd@udit.es}{dpd@udit.es} ·
        \href{mailto:proteccion.datos@udit.es}{proteccion.datos@udit.es}
\end{itemize}

\section*{8.\ Comité de ética}

Este protocolo se somete al CEI de UDIT
(\href{mailto:comite.etica.investigacion@udit.es}{comite.etica.investigacion@udit.es}).
Hasta su autorización, \textbf{ningún} artefacto vuestro se usa
como dato de investigación publicable.

\newpage
\LogoRow

\begin{center}
  {\Large\sffamily\bfseries Formulario de consentimiento informado}\\[0.35em]
  {\normalsize\sffamily \TTODfull\ / Front-End~II · versión 2026-09-26}
\end{center}

\vspace{0.4em}
\noindent
He leído la \textbf{Hoja de información a participantes} (versión 2026-09-26)
que precede a este formulario y he podido hacer preguntas.

\vspace{0.5em}
\noindent
Marque con claridad cada casilla que acepte
(\texttt{$\square$}~=~no marcado; escriba \textbf{Sí} o marque a mano):

\vspace{0.4em}
\renewcommand{\arraystretch}{1.35}
\begin{tabularx}{\textwidth}{@{}c>{\raggedright\arraybackslash}X c@{}}
  \toprule
  \textbf{\#} & \textbf{Declaración} & \textbf{Sí}\\
  \midrule
  \textbf{A} &
  Comprendo la actividad \textbf{docente} del grupo
  \TTODfull\
  (participar como desarrollador/a en el proyecto de la asignatura)
  y acepto participar en ella según la guía docente.
  & $\square$ \\
  \textbf{B} &
  Consiento el \textbf{uso investigador secundario y seudonimizado}
  de mis artefactos de proceso ya producidos para la evaluación
  (commits/PRs, declaración de uso de IA, notas estructuradas de defensa),
  \textbf{solo} tras autorización ética y sin efecto sobre mi calificación
  si me niego o me retiro.
  & $\square$ \\
  \textbf{C} &
  Comprendo que, en el corpus investigador, \textbf{no} se recogen categorías
  especiales de datos ni se publicarán mis identificadores directos.
  & $\square$ \\
  \textbf{D} &
  \emph{(Opcional --- solo si se activa en el futuro)}
  Consiento encuesta~/ entrevista adicional según información específica
  que se me entregaría entonces.
  & N/A \\
  \textbf{E} &
  \emph{(Opcional --- solo si se activa)}
  Consiento grabación de audio/vídeo de defensa o entrevista según
  información específica.
  & N/A \\
  \textbf{F} &
  \emph{(Opcional)}
  Consiento la difusión de extractos de código o capturas
  \textbf{no identificativas} en publicaciones académicas
  (consentimiento por artefacto si se me pide después).
  & $\square$ \\
  \textbf{G} &
  \emph{(Opcional)}
  Autorizo el uso de \textbf{mi imagen} (fotografías / capturas del
  grupo) por UDIT y el proyecto \TTODfull\ en comunicaciones
  relacionadas: LinkedIn y otras redes, medios, publicaciones,
  revistas internas universitarias y sitio web público, incluido el
  material que gestione el departamento de comunicación.
  & $\square$ \\
  \bottomrule
\end{tabularx}

\vspace{1.2em}
\noindent
\textbf{Nombre y apellidos:}
\hrulefill

\vspace{1.1em}
\noindent
\begin{tabularx}{\textwidth}{@{}X X@{}}
  \textbf{Firma:} \hrulefill &
  \textbf{Fecha:} \underline{\hspace{1.2cm}}~/~\underline{\hspace{1.2cm}}~/~\underline{\hspace{2.2cm}}
\end{tabularx}

\vspace{1.4em}
\noindent{\color{rulegray}\hrule height 0.4pt}
\vspace{0.5em}

\noindent{\small\color{uditink}
\textbf{Uso interno} (no completar con datos de terceros en la copia del estudiante):
código de seudónimo \texttt{P\_\_\_\_} · custodio: \underline{\hspace{3.5cm}} ·
archivo fuera de git (ver protocolo de custodia).
Las copias \textbf{firmadas} no se suben al repositorio del proyecto TTOD.}

\vspace{0.8em}
\noindent{\footnotesize\color{uditnavy}
Grupo ECSIT · Web: \EcsitLink.}

\vspace{0.35em}
\LegalEntities

\end{document}
